\include{zLaTex-Preamble}

\title{Project 3 - Design and Test Writeup} 
\author{Josh Miller}
%% \date{}

\begin{document}
\rhead{\fancyplain{}{Project 3 Design - Josh Miller}}
\maketitle
\lang{C}

\section{Purpose}

The purpose of this project is to implement different locking
mechanisms for the purpose of testing the overhead introduced in
resource contention.  

\section{Design}

\subsection{Framework}

The overhead introduced by each lock will first be tested by having
{\tt N} threads attempting to update a single incrementing counter.
Because this setup is serliazed completely by the single shared
resource, I expect it to be impossible that any locked/threaded setup
performs with higher performance than a single threaded, non-locked
loop.  Therefore, we will be comparing the speedup of the threaded
versions versus the serial version and expect the speedup, $S$, to be
$0 < S < 1$.

The framework for testing these locks take a single {\tt Lock}
structure.  This structure contains a {\tt Union} of the five lock
implementations that will be implement for this project.  As the
threads are spawned, they are passed function pointers which specify
the procedures for the {\tt acquire} and {\tt release} of each
lock. This allows for the simple and efficient abstraction of the lock
over the thread loop while fixing each lock and unlock call to a
function pointer de-reference. This may however increase the over head
by disallowing the compiler to in-line the code for the lock and
unlock, resulting in greater branch overhead.

\begin{code}
while (LOOP_THREADS){
        args->acquire(args->lock);
#ifdef TEST
        args->test_counter ++;
#endif
        *args->counter += 1;
        args->release(args->lock);
}
\end{code}

The global variable {\tt LOOP\_THREADS} is switched off after a given
duration of time by the main thread.  Timing of the functions is not
necessary because we are controlling the duration of execution, rather
than the number of iterations.  The duration which to sleep is
implemented using the {\tt usleep()} method.

The second round if testing will use the value of the counter as the
loop test condition. The dependant variable to be tested will then be
the time required to complete the execution of the loop a given number
of times.

Included in the thread loop is a preprocessor switch which allows the
thread loop to maintain its own counter during testing.  After threads
are joined, the counters are totalled and compared against the shared
counter.  In this way, I can test the accuracy of each of the locks.
Any locks which do not provide mutual exclusion will allow race
conditions, and the final counter comparison will fail. Therefore, I
was able to write one test to test all of the locks.  By calling {\tt
  thread\_battery()} I am able to reinitialize the lock based on the
implementation ad pass the {\tt *Lock} and its {\tt acquire} and {\tt
  release} functions to the test.

\begin{code}
Lock lock;

TASB_init(&lock);
thread_battery("TASB", &lock, &TASB_lock, &TASB_unlock);

ALOCK_init(&lock, 8);
thread_battery("ALOCK", &lock, &ALOCK_lock, &ALOCK_unlock);
\end{code}

An example output of the test run on the {\tt ALOCK\_t} lock
implementation is as follows.

\lang{test}
\begin{code}
TEST: [               ALOCK] [                 1 threads 1 trials] [ PASSED ]
TEST: [               ALOCK] [                 2 threads 1 trials] [ PASSED ]
TEST: [               ALOCK] [                 4 threads 1 trials] [ PASSED ]
TEST: [               ALOCK] [                 8 threads 1 trials] [ PASSED ]
TEST: [               ALOCK] [                32 threads 1 trials] [ PASSED ]
TEST: [               ALOCK] [                 8 threads 5 trials] [ PASSED ]
TEST: [               ALOCK] [                8 threads 10 trials] [ PASSED ]
TEST: [               ALOCK] [                8 threads 23 trials] [ PASSED ]
\end{code}

Futhermore, each lock will have a {\tt TryLock} method associated with
it which will return {\tt False} if it was unable to acquire the lock
and {\tt True} if the lock was aqcuired.

\lang{C}
\subsection{TASLock}

The following code is my implementation for the {\tt TASLock}

\begin{code}
int TAS_lock(Lock *lock)
{
    while (__sync_lock_test_and_set(&lock->tas, 1));
    return 0;
}

int TAS_unlock(Lock *lock)
{
    return !__sync_lock_test_and_set(&lock->tas, 0);
}

int TAS_init(Lock *lock)
{
    return (lock->tas = 0);
}
\end{code}

This uses the built-in intel atomic commands to acquire the lock.


\section{TASLock Backoff}

The following code is my implementation for the {\tt TASBLock}

\begin{code}
int TASB_init(Lock *lock)
{
    srand(time(NULL));
    return (lock->tas = 0);
}

int TASB_lock(Lock *lock)
{
    int delay = 1;
    int delay_top = 1;
    while (__sync_lock_test_and_set(&lock->tas, 1)){
        delay_top <<= 1;
        delay = rand()/(RAND_SCALE/delay_top);
        usleep(delay);
    }
    return 0;
}

int TASB_unlock(Lock *lock)
{
    return !__sync_lock_test_and_set(&lock->tas, 0);
}
\end{code}

Every time {\tt TASB\_lock} is called, backoff is initially capped at
1 microsecond.  While the thread is unable to acquire the lock, it
will double the maximum delay to {\tt delay\_top} microseconds.  After
failure to acquire the lock, the thread will then wait on a delay in
the range of 0 to {\tt delay\_top} microseconds.

\subsection{pthread\_mutex\_t}

The pthread is simply included in the lock {\tt Union} and is
initialized, locked, and unlocked in the normal way.

\subsection{ALOCK}



\begin{code}
#define BLOCKED 1
#define FREE 0
#define PAD 8

int ALOCK_init(Lock *lock, int n)
{
    lock->alock.n = n*PAD;
    lock->alock.tail = 0;
    lock->alock.curr = 0;
    lock->alock.flag = (volatile int*) malloc(lock->alock.n*sizeof(int));
    for (int i = 1; i < lock->alock.n; i++)
        lock->alock.flag[i] = BLOCKED;
    lock->alock.flag[0] = FREE;
    return 0;
}

int ALOCK_free(Lock *lock)
{
    free(lock->alock.flag);
    return 0;
}


int ALOCK_lock(Lock *lock)
{

    int i = __sync_fetch_and_add(&lock->alock.tail, PAD) % lock->alock.n;
    while (lock->alock.flag[i] == BLOCKED);
    return 0;

}

int ALOCK_unlock(Lock *lock)
{

    int i = lock->alock.curr % lock->alock.n;
    lock->alock.flag[i] = BLOCKED;
    lock->alock.curr += PAD;
    i = lock->alock.curr % lock->alock.n;
    lock->alock.flag[i%lock->alock.n] = FREE;
    
    return 0;
}
\end{code}

The Anderson lock was implemented with a padding of 8 bytes.

\subsection{CLH Queue Lock}

As of this writeup, I have not implemented the CLH Queue lock. The
design will follow that described in the book.  Because C is not a
garbage collecting language, the initializer function will initialize
the total number, $N$, nodes needed for $N$ threads. Unlocking will
not involve freeing the node, but rather allowing it to be acquired by
another thread.  The principle of this lock is that each thread waits
for its predecessor to claim that it has released the lock.

\subsection{The awesomest lock ever to be devised}

I will come up with this lock after I have run performance tests on
the above locking methods and identified the factor with the largest
performance impact.  I will build the lock with an attempt to address
this factor.

\section{Packet Processing}

This project will require the following updates to Project 2:

\begin{itemize}
\item the packet processing loop test will be based on ellapsed time,
  rather than number of packets processed
\item each thread will not have a specific source/destination
  immutably assigned to it. Therefore each queue will have a lock
  assigned to it that a thread must acquire before processing a packet
  in the queue.
\end{itemize}

I believe these design changes are self-explainatory and will not
detail what I believe is an obvious implementation.

\emph{I will be able to reuse the testing framework for Project 2 to
  assess the accuracy of these implementations}

\section{Experiments}

As described in the assignment, these are the experiments to be
carried out with my hypothesis for each of them.

\subsection{Counter Tests}
\begin{enumerate}

\item 
\textbf{ Idle Lock Overhead 1} Run SerialCounter for the Time-based
counter experiment to find the throughput (increments / ms) of a
single processor incrementing a counter (a volatile, which admittedly
has some overhead). Then, run ParallelCounter with n = 1 and all L ($\in$ 
\{TASLock, BackoffLock, Mutex, ALock, CLHLock, MCSLock\}). 

\emph{Hypothesis:} The SerialCounter version will have the best
performance. Therefore, the speedup will be $S \leq 1$ for all locks.
Furthermore, I predect the following relations between locks:
\begin{itemize}
\item The TASLOCK will be faster that the BackoffLock. I predict this
  because I believe that the delay introduced by the contention over
  the shared resource will be less than the delay artificially
  introduced by the backoff.
\item I am unsure of the releative performance of the {\tt
  pthread\_mutex\_t} lock because I am unfamiliar with its
  implementation
\item The ALOCK will execute slower than TestAndSet locks.  This is
  because the lock contention is distributed over an array, rather
  then centralized on one integer. However, there is no contention
  with a single thread in the threaded version.
\item the CLHLock will perform worse than the ALOCK.  I am not
  confident in this prediction, but I'm making it because I believe
  the implementation of the CLHLock will require more overhead to
  manage the linked list queue due to its complexity over an array
  queue.
\end{itemize}


\item \textbf{Idle Lock Overhead 2} Run SerialCounter for the Work-based
  experiment to find the throughput (increments / ms) of a single
  processor incrementing a counter (a volatile, which admittedly has
  some overhead). Then, run ParallelCounter with n = 1 and all L (∈
  {TASLock, BackoffLock, Mutex, ALock, CLHLock, MCSLock}). 

  Provide the speedup (i.e. ratio of ParallelCounter throughput to
  SerialCounter) for each L in a table. Do these results make sense,
  given the relative complexity of the uncontended path through each
  lock method? How do the results for the two different (time vs work
  based) approaches compare?  (Depending on your implementations there
  may not be differences). You are responsible for justi- fying the
  answer to these questions. Your answers should persuade graders that
  you understand the internals of your code. If you cannot explain a
  result with data (taken from experiments or from code analysis) then
  you should ask for help.

\emph{Hypothesis:} Because both this and the previous test are
functionally indesdinguishable (as far as I can tell) tests of
throughput, the performance differences between implementations will
be the same as in Experiment~1.

\item \textbf{Lock Scaling 1} Optimize the DELAY times for ParallelCounter with
  L = BackoffLock for n = 8. Run ParallelCounter for all L and n ∈ {1,
    2, 4, 8, 16, 32, 64} with the tuned version of BackoffLock. 

  Using the time-based implementation, plot speedup (ratio of
  ParallelCounter throughput to SerialCounter - it should be less than
  1) for each L (i.e. one curve for each L) on the Y -axis vs. n on
  the X-axis. Describe on your results - what is your hypothesis about
  the performance, particularly of the more complex locking
  algorithms? Back up this answer with data that demonstrates your
  understanding.


\emph{Hypothesis:} The SerialCounter version will have the best
performance. Therefore, the speedup will be $S \leq 1$ for all locks.
Furthermore, I predect the following relations between locks:

\item\textbf{ Lock Scaling 2} Optimize the DELAY times for
  ParallelCounter with L = BackoffLock for n = 8. Run the work-based
  approach to ParallelCounter for all L and n $\in$ {1, 2, 4, 8, 16,
    32, 64} with the tuned version of BackoffLock.

\emph{Hypothesis:} I have no prediction for the actual delay
configuration settings that will be optimal, but I predict that there
is such a configuration that the exponential backoff lock will
outperform the simple test and set lock.

\item\textbf{ Fairness }Find the standard deviation of the count for each
  worker in the time-based Lock Scaling experiment for n = C, where C
  is the number of cores on your test machine. Produce a scatter plot
  of throughput (Y -axis) vs. standard deviation of counts among
  workers (X-axis) with each lock, L, represented. 

\emph{Hypothesis:} I do believe it will corroborate my hypothesis from 
the time-based lock experiment.

\end{enumerate}

\subsection{Packet Tests}
\begin{enumerate}


\item\textbf{ Idle Lock Overhead} Run ParallelPacket with n = 1, S
  $\in$ \{LockFree, HomeQueue\}, all L and W $\in$ \{25, 50, 100, 200,
  400, 800\} on the uniformly distributed packets. 

Plot the speedup of
  ParallelPacket (with S = HomeQueue) throughput relative to
  ParallelPacket (with S = LockFree) for each L on the Y -axis vs. W
  on the X-axis. 

\emph{Hypothesis:} I predict that the speedup of the ParallelPacket
with (S = HomeQueue) will perform worse than the runs with (S =
LockFree).  This is because the tests are identical save additional
overhead.


\item\textbf{ Speedup with Uniform Load} Run ParallelPacket and SerialPacket on
  the uniformly distributed packets with the number of workers n $\in$ \{1,
  2, 3, 7, 15\}, W $\in$  \{1000, 2000, 4000, 8000\}, S $\in$  \{LockFree,
  RandomQueue, LastQueue\} and L $\in$  \{TASLock, BackoffLock, Mutex,
    ALock, CLHLock, MCSLock\}.  

\emph{Hypothesis:} There are many things happening in this
experiment. Given this large search space, my hypothesis will be
slightly vague:
\begin{itemize}
\item I predict that all parallel versions will outperform the serial
  version when the number of worker threads is $\geq$ 2. This means
  that I believe the overhead introduced per packed by any of the
  locks implemented will not be equal to the cost of processing a
  single packet in serial.
\item I predict that the performance of the parallel versions with
  load balancing and locking will be less than the fixed queue
  threading scheme. This is because the unbalance in the load is small
  enough that the overhead from the locks is not outweighed by the
  balancing sheme.
\item I predict that speedup of each of the threaded version will
  asymptotically approach 15 for 15 worker threads as the mean packet
  work approaches 8000
\end{itemize}


\item\textbf{ Speedup with Exponential Load} Repeat the previous experiment,
  except with the exponentially distributed packets. 

\emph{Hypothesis:} I maintain all predicions from the previous
experiment save the one about the comparison between locking thread
schemes and lock free thread schemes.  I believe that the locked
threads will outperform the unlocked threads because the overhead of
the locking is outweighed by the benefit in load balancing.
Furthermore, I predict that the discrepency between the locked and
unlocked thread schemes will increase as the mean packet length
increases. This is because the load imbalance is mor noticable for
larger loads.

\item\textbf{ Speedup with Awesome} Performance testing shemes in my
  personal design have yet to be devised


\end{enumerate}




\end{document}
